\section{Examples}
\label{section:examples}

This section walks through typical Proteo workflows with JSON configurations.

\subsection{Minimal single run}

Use the bundled test configuration \texttt{Codes/test.json} (three phases, two resizes, multiple stage types including I/O):

\begin{lstlisting}[style=bashstyle,caption={Run test.json with Slurm auto-allocation.},captionpos=b]
cd Results
bash ../Exec/singleRun.sh ../Codes/test.json mypartition 0 1 0
bash ../Exec/CheckRun.sh test 0 1 4 3 100
\end{lstlisting}

Expected outputs (index~0): \texttt{R0\_Global.out}, \texttt{R0\_G0NP2.out}, \texttt{R0\_G1NP4.out}, \texttt{R0\_G2NP6.out}, plus \texttt{slurm-*.out}.

\subsection{Custom allocation}

Obtain a Slurm allocation first, then run inside it:

\begin{lstlisting}[style=bashstyle,caption={Run inside an sbatch allocation.},captionpos=b]
sbatch -N 2 -p mypartition --wrap \
  "bash Exec/singleRunCostum.sh Codes/test.json 0 5 0"
\end{lstlisting}

This runs five repetitions with output prefix \texttt{R0}. Use a different \texttt{outIndex} (2nd argument) to distinguish batches: \texttt{R5\_*}, \texttt{R6\_*}, etc.

\subsection{Creating a config with GenConfig}

\begin{enumerate}
  \item Start GenConfig (\texttt{python3 app.py} in \texttt{GenConfig/}).
  \item Stay in \textbf{Single} mode. Click \textbf{New} to load the \texttt{test.json} template.
  \item Adjust \texttt{Procs} and \texttt{Iters} in the groups section.
  \item Click \textbf{Validate}, then \textbf{Download JSON} as \texttt{my\_experiment.json}.
  \item Inside an allocation: \texttt{bash Exec/singleRunCostum.sh my\_experiment.json 0 1 0}.
\end{enumerate}

\subsection{Simple resize example}

Listing~\ref{code:json-full-example} (Section~\ref{section:json-config}) defines one phase with Monte Carlo + Allreduce + Allgather, resizing from 2 to 10 processes after one source-group iteration. Save it as \texttt{simple\_resize.json} and run:

\begin{lstlisting}[style=bashstyle]
bash Exec/singleRunCostum.sh simple_resize.json 0 1 0
\end{lstlisting}

\subsection{Multi-phase application}

Phase~0 of \texttt{Codes/test.json} mixes I/O read, Bcast, and compute:

\begin{lstlisting}[language=json,caption={Phase 0 of Codes/test.json.},captionpos=b,label=code:json-phase-excerpt]
{
  "Total_Iters": 1,
  "Total_Stages": 3,
  "stages": [
    {
      "Stage_Type": 10,
      "Stage_Bytes": 0,
      "Stage_Involved_Procs": 2,
      "Stage_Time": 4,
      "Granularity": 1000
    },
    { "Stage_Type": 5, "Stage_Bytes": 1000, "Stage_Time": 0 },
    {
      "Stage_Type": 0,
      "Stage_Bytes": 0,
      "Stage_Time": 1,
      "Granularity": 10000
    }
  ]
}
\end{lstlisting}

Phase~1 runs 22 iterations with compute + two Allreduce stages + Allgather; phase~2 performs a single large I/O write. This illustrates how phases chain different stage sequences within the same process group.

\subsection{Combinatorial study}

Variant fields in \texttt{GenConfig/examples/complex.example.json}:

\begin{lstlisting}[language=json,caption={Variant fields in complex.example.json.},captionpos=b]
"SDR": "500000:1000000",
"ADR": "0:50",
"Procs": "2:4",
"FactorS": "1:1"
\end{lstlisting}

\textbf{Workflow:}
\begin{enumerate}
  \item Edit variants in GenConfig Multi mode (or edit the file directly).
  \item Validate to see the predicted number of valid combinations.
  \item Download as \texttt{complex.json}.
  \item Expand: \texttt{python3 Exec/PythonCodes/read\_multiple\_json.py complex.json sweep\_}.
  \item Run each \texttt{Desglosed-*/sweep\_N.json} with \texttt{singleRun.sh}, varying partition and output index as needed.
\end{enumerate}

\subsection{Conjugate-gradient-like profile}

A CG-like iteration can be approximated with one compute stage plus several communication stages (see Listing~\ref{code:json-full-example} for structure). Use matrix--vector compute (\texttt{Stage\_Type}: 1), small Allreduce messages (\texttt{Stage\_Bytes}: 8), and a large Allgather (\texttt{Stage\_Bytes} set explicitly). Tune \texttt{Granularity} and \texttt{Stage\_Time} to match target iteration time on a reference process count, then set \texttt{FactorS} per group for scaling behaviour.

\subsection{Error checking after a batch}

After multiple runs sharing prefix \texttt{exp}:

\begin{lstlisting}[style=bashstyle]
bash Exec/CheckRun.sh exp 9 5 4 3 120
\end{lstlisting}

Checks indices~0 through~9, five repetitions each, expecting four stages and three groups, rejecting iterations longer than 120~seconds.
